\documentclass[11pt]{letter}
\usepackage[margin=1in]{geometry}
\usepackage{parskip}
\usepackage{hyperref}

\signature{Ning Zhou\\
Public R\&D Center, SuperMap Software Co., Ltd.\\
Beijing, China\\
\href{mailto:zhouning1@supermap.com}{zhouning1@supermap.com}}

\address{}

\begin{document}

\begin{letter}{The Editor \\ \textit{Nature Sustainability}}

\opening{Dear Editor,}

We are pleased to submit our manuscript, ``Model-based AI planning enables county-scale farmland consolidation in fragmented mountain landscapes,'' for consideration as a Research Article in \textit{Nature Sustainability}.

Mountainous and hilly croplands account for a substantial share of global food production, yet face an acute productivity bottleneck: extreme parcel fragmentation prevents agricultural mechanisation, increases labour costs, and threatens food-security gains under climate stress. Land consolidation, the spatial rearrangement of farmland and forest parcels under regulatory and topographic constraints, is the policy instrument used to address this bottleneck. Generating defensible consolidation plans at county scale, however, requires evaluating combinatorial trade-offs across tens of thousands of cadastral parcels under hard regulatory constraints. Manual planning cannot scale to this size, and off-the-shelf reinforcement learning either fails to converge or requires eight to twelve GPU-hours per scenario with non-trivial collapse rates, making the iterative scenario exploration that planners actually need impractical.

Our manuscript makes four connected contributions that we believe align with the journal's interest in actionable sustainability science.

First, the diagnostic. Standard mean-squared-error training of a learned planning model masks a hidden ranking failure: on real cadastral data from Bishan District, Chongqing ($52{,}515$ parcels), per-action reward predictions are essentially random (pairwise accuracy $51.6\%$) despite cosine similarity above $0.998$. Four independent operational failures (dream-policy collapse, MPC ceiling at $-0.209\%$ slope, greedy-continuation reversal, retrained-baseline persistence) are explained by this single gap.

Second, the fix. Adding an auxiliary pairwise large-margin loss to the existing reward head raises ranking accuracy to $85.5\%$ without architectural change, and converts model-predictive control from failure to the best deterministic-exploitation strategy we tested. The same intervention reproduces on a second county (Neijiang Dongxing District, $76{,}376$ parcels) and across seven synthetic stress-test landscapes spanning terrain, size, and fragmentation regimes. The cross-county and cross-landscape replication is essential to the policy claim: the method is not tuned to a single morphology.

Third, the deployment. We release a four-tool ArcGIS~Pro Python Toolbox that takes a county DLTB shapefile and digital elevation model and returns an optimised cadastral plan, with a 70-second smoke test for continuous integration and approximately $7$~min per planning scenario on a desktop CPU. Together with the open synthetic benchmark (CC-BY 4.0; deterministic generator), this allows planners to explore consolidation scenarios in minutes rather than months, and lets external researchers benchmark future methods without access to sensitive cadastral data.

Fourth, the verification. We anticipate that reviewers will ask whether the headline $-2.00\%$ slope reduction is real geometry or a statistical artefact of the learned ensemble. We address this with a five-layer verification stack disclosed openly in the manuscript and reproducible by any reader from the released artefacts: (i)~independent GIS recomputation of all three primary metrics directly from the optimised cadastral shapefile, agreeing with the planner-reported numbers to floating-point noise without invoking any learned component; (ii)~ensemble-member subset ablation (n=3, all 2-of-3 subsets) replacing the deterministic-pipeline std=0 with a real $\pm 0.026$\,pp slope variance, with the full-ensemble result sitting at the median of the subset distribution; (iii)~direct measurement of the ensemble's 1-step prediction error against the true environment along the same 100-step trajectory MPC visits, finding the ensemble strongly mis-calibrated as a regressor (reward MAE $0.82$) yet preserving Spearman rank correlation $+0.22$ that is operationally sufficient at H$=$5; (iv)~a counter-factual ablation that replaces the ensemble entirely with the perfect dynamics model and finds it strictly worse on all three objectives, proving the ensemble's contribution is structural depth/breadth at near-constant cost rather than statistical accuracy; and (v)~cross-county replication on Neijiang. To our knowledge this is the deepest verification stack offered in the published RL-for-land-use literature.

The manuscript was not submitted, in whole or in part, to any other journal. The authors have no competing interests. We confirm that all data presented are original, that no identifying cadastral records are redistributed, and that the released benchmark contains only synthetic landscapes derived from a deterministic seed-controlled generator. We thank the Editor and reviewers for their time and look forward to their assessment.

\closing{Sincerely,}

\end{letter}
\end{document}
